
\begin{table*}[!th]
    \centering
    \small
    \caption{
        \textbf{Prompt used to check \leetcode questions.} {\color{blue} Blue text} denote instance-specific inputs. {\color{violet} Violet text} denotes the output from GPT-4.
        For each question, we provide the title and the problem statement.
        We use six in-context demonstrations, and the instruction prompts the model to categorize the question based on the criteria provided.
        Although there are three possible labels, we only keep questions that were labeled as 2, which are grounded in real-world concepts.
        For sake of brevity, we only show one example, but other examples can be found on the code repo. 
        % For each question, we provide the original question and answer, as well as the theorem name and theorem definition from the original TheoremQA dataset.
        % The instruction prompts the model to rewrite the question with a different surface form without changing the reasoning steps or the final answer.
        % We hand-write two examples to illustrate the rewriting process.
        % We also allow the model to skip the question if it is too difficult to rewrite.
    }
    \label{tab:leetcode_check}
    % \resizebox{0.98\linewidth}{!}{
    \begin{tabular}{>{\raggedright\arraybackslash\tt}p{0.98\textwidth}<{}}
        \toprule
            \vspace{-1em}
            Read the coding question and categorize it using the following criteria:\newline%
 0: The question is not grounded in any real{-}world concepts. The description only uses coding{-}specific terms, such as "linked list", "binary search", "palindrome", "sorting", etc..\newline%
 1: The question is not grounded in any real{-}world concepts or real{-}world concepts that are commonly used in the context of coding, such as needle in a haystack, strings/words, or a spiral matrix.\newline%
 2: The question is grounded in real{-}world concepts that are not commonly used in the context of coding, such as building height, planting trees, or games. It may still uses some code{-}specific terms to specify the data structure involved.\newline%
\newline%
You should only consider the initial problem statement and problem title, not the examples or constraints.\newline%
Use the following examples to help you categorize the question:\newline%
\newline%
Example 1:\newline%
\{\{\newline%
~~~~"title": "Merge Two Sorted Lists",\newline%
~~~~"question": "You are given the heads of two sorted linked lists `list1` and `list2`.\newline%
Merge the two lists in a one **sorted** list. The list should be made by splicing together the nodes of the first two lists.\newline%
\newline%
Return \_the head of the merged linked list\_.\newline%
\newline%
**Example 1:**\newline%
% \newline%
\textit{rest of the example omitted...}
\newline%
% **Input:** list1 = \textbackslash{}{[}1,2,4\textbackslash{}{]}, list2 = \textbackslash{}{[}1,3,4\textbackslash{}{]}\newline%
% **Output:** \textbackslash{}{[}1,1,2,3,4,4\textbackslash{}{]}\newline%
% \newline%
% **Example 2:**\newline%
% \newline%
% **Input:** list1 = \textbackslash{}{[}\textbackslash{}{]}, list2 = \textbackslash{}{[}\textbackslash{}{]}\newline%
% **Output:** \textbackslash{}{[}\textbackslash{}{]}\newline%
% \newline%
% **Example 3:**\newline%
% \newline%
% **Input:** list1 = \textbackslash{}{[}\textbackslash{}{]}, list2 = \textbackslash{}{[}0\textbackslash{}{]}\newline%
% **Output:** \textbackslash{}{[}0\textbackslash{}{]}\newline%
% \newline%
% **Constraints:**\newline%
% \newline%
% *   The number of nodes in both lists is in the range `{[}0, 50{]}`.\newline%
% *   `{-}100 <= Node.val <= 100`\newline%
% *   Both `list1` and `list2` are sorted in **non{-}decreasing** order."\newline%
\}\}\newline%
\newline%
\{\{\newline%
    "label": 0\newline%
\}\}\newline%
\newline%
\textit{Examples 2-6 are omitted for brevity...}\newline%
\newline%
% Example 2:\newline%
% \{\{\newline%
%     "title": "Find the Index of the First Occurrence of a String",\newline%
%     "question": "Given two strings `needle` and `haystack`, return the index of the first occurrence of `needle` in `haystack`, or `{-}1` if `needle` is not part of `haystack`.\newline%
% \newline%
% **Example 1:**\newline%
% \textit{rest of the example omitted...}
% \newline%
% **Input:** haystack =  "sadbutsad ", needle =  "sad "\newline%
% **Output:** 0\newline%
% **Explanation:**  "sad " occurs at index 0 and 6.\newline%
% The first occurrence is at index 0, so we return 0.\newline%
% \newline%
% **Example 2:**\newline%
% \newline%
% **Input:** haystack =  "leetcode ", needle =  "leeto "\newline%
% **Output:** {-}1\newline%
% **Explanation:**  "leeto " did not occur in  "leetcode ", so we return {-}1.\newline%
% \newline%
% **Constraints:**\newline%
% \newline%
% *   `1 <= haystack.length, needle.length <= 104`\newline%
% *   `haystack` and `needle` consist of only lowercase English characters."\newline%
% \}\}\newline%
% \newline%
% \{\{\newline%
%     "label": 1\newline%
% \}\}\newline%
% \newline%
% Example 3:\newline%
% \{\{\newline%
%     "title": "Dota2 Senate",\newline%
%     "question": "In the world of Dota2, there are two parties: the Radiant and the Dire.\newline%
% \newline%
% The Dota2 senate consists of senators coming from two parties. Now the Senate wants to decide on a change in the Dota2 game. The voting for this change is a round{-}based procedure. In each round, each senator can exercise **one** of the two rights:\newline%
% \newline%
% *   **Ban one senator's right:** A senator can make another senator lose all his rights in this and all the following rounds.\newline%
% *   **Announce the victory:** If this senator found the senators who still have rights to vote are all from the same party, he can announce the victory and decide on the change in the game.\newline%
% \newline%
% Given a string `senate` representing each senator's party belonging. The character `'R'` and `'D'` represent the Radiant party and the Dire party. Then if there are `n` senators, the size of the given string will be `n`.\newline%
% \newline%
% The round{-}based procedure starts from the first senator to the last senator in the given order. This procedure will last until the end of voting. All the senators who have lost their rights will be skipped during the procedure.\newline%
% \newline%
% Suppose every senator is smart enough and will play the best strategy for his own party. Predict which party will finally announce the victory and change the Dota2 game. The output should be `"Radiant "` or `"Dire "`.\newline%
% \newline%
% **Example 1:**\newline%
% \newline%
% **Input:** senate =  "RD "\newline%
% **Output:**  "Radiant "\newline%
% **Explanation:** \newline%
% The first senator comes from Radiant and he can just ban the next senator's right in round 1. \newline%
% And the second senator can't exercise any rights anymore since his right has been banned. \newline%
% And in round 2, the first senator can just announce the victory since he is the only guy in the senate who can vote.\newline%
% \newline%
% **Example 2:**\newline%
% \newline%
% **Input:** senate =  "RDD "\newline%
% **Output:**  "Dire "\newline%
% **Explanation:** \newline%
% The first senator comes from Radiant and he can just ban the next senator's right in round 1. \newline%
% And the second senator can't exercise any rights anymore since his right has been banned. \newline%
% And the third senator comes from Dire and he can ban the first senator's right in round 1. \newline%
% And in round 2, the third senator can just announce the victory since he is the only guy in the senate who can vote.\newline%
% \newline%
% **Constraints:**\newline%
% \newline%
% *   `n == senate.length`\newline%
% *   `1 <= n <= 104`\newline%
% *   `senate{[}i{]}` is either `'R'` or `'D'`."\newline%
% \}\}\newline%
% \newline%
% \{\{\newline%
%     "label": 2\newline%
% \}\}\newline%
% \newline%
% Example 4:\newline%
% \{\{\newline%
%     "title": "Shortest Palindrome",\newline%
%     "question": 'You are given a string `s`. You can convert `s` to a palindrome by adding characters in front of it.\newline%
% \newline%
% Return \_the shortest palindrome you can find by performing this transformation\_.\newline%
% \newline%
% **Example 1:**\newline%
% \newline%
% **Input:** s = "aacecaaa"\newline%
% **Output:** "aaacecaaa"\newline%
% \newline%
% **Example 2:**\newline%
% \newline%
% **Input:** s = "abcd"\newline%
% **Output:** "dcbabcd"\newline%
% \newline%
% **Constraints:**\newline%
% \newline%
% *   `0 <= s.length <= 5 * 104`\newline%
% *   `s` consists of lowercase English letters only.'\newline%
% \}\}\newline%
% \newline%
% \{\{\newline%
%     "label": 0\newline%
% \}\}\newline%
% \newline%
% Example 5:\newline%
% \{\{\newline%
%     "title": "Course Schedule II",\newline%
%     "question": 'There are a total of `numCourses` courses you have to take, labeled from `0` to `numCourses {-} 1`. You are given an array `prerequisites` where `prerequisites{[}i{]} = {[}ai, bi{]}` indicates that you **must** take course `bi` first if you want to take course `ai`.\newline%
% \newline%
% *   For example, the pair `{[}0, 1{]}`, indicates that to take course `0` you have to first take course `1`.\newline%
% \newline%
% Return \_the ordering of courses you should take to finish all courses\_. If there are many valid answers, return **any** of them. If it is impossible to finish all courses, return **an empty array**.\newline%
% \newline%
% **Example 1:**\newline%
% \newline%
% **Input:** numCourses = 2, prerequisites = \textbackslash{}{[}\textbackslash{}{[}1,0\textbackslash{}{]}\textbackslash{}{]}\newline%
% **Output:** \textbackslash{}{[}0,1\textbackslash{}{]}\newline%
% **Explanation:** There are a total of 2 courses to take. To take course 1 you should have finished course 0. So the correct course order is \textbackslash{}{[}0,1\textbackslash{}{]}.\newline%
% \newline%
% **Example 2:**\newline%
% \newline%
% **Input:** numCourses = 4, prerequisites = \textbackslash{}{[}\textbackslash{}{[}1,0\textbackslash{}{]},\textbackslash{}{[}2,0\textbackslash{}{]},\textbackslash{}{[}3,1\textbackslash{}{]},\textbackslash{}{[}3,2\textbackslash{}{]}\textbackslash{}{]}\newline%
% **Output:** \textbackslash{}{[}0,2,1,3\textbackslash{}{]}\newline%
% **Explanation:** There are a total of 4 courses to take. To take course 3 you should have finished both courses 1 and 2. Both courses 1 and 2 should be taken after you finished course 0.\newline%
% So one correct course order is \textbackslash{}{[}0,1,2,3\textbackslash{}{]}. Another correct ordering is \textbackslash{}{[}0,2,1,3\textbackslash{}{]}.\newline%
% \newline%
% **Example 3:**\newline%
% \newline%
% **Input:** numCourses = 1, prerequisites = \textbackslash{}{[}\textbackslash{}{]}\newline%
% **Output:** \textbackslash{}{[}0\textbackslash{}{]}\newline%
% \newline%
% **Constraints:**\newline%
% \newline%
% *   `1 <= numCourses <= 2000`\newline%
% *   `0 <= prerequisites.length <= numCourses * (numCourses {-} 1)`\newline%
% *   `prerequisites{[}i{]}.length == 2`\newline%
% *   `0 <= ai, bi < numCourses`\newline%
% *   `ai != bi`\newline%
% *   All the pairs `{[}ai, bi{]}` are **distinct**.'\newline%
% \}\}\newline%
% \newline%
% \{\{\newline%
%     "label": 2\newline%
% \}\}\newline%
% \newline%
% Example 6:\newline%
% \{\{\newline%
%     "title": "Concatenated Words",\newline%
%     "question": 'Given an array of strings `words` (**without duplicates**), return \_all the **concatenated words** in the given list of\_ `words`.\newline%
% \newline%
% A **concatenated word** is defined as a string that is comprised entirely of at least two shorter words (not necesssarily distinct) in the given array.\newline%
% \newline%
% **Example 1:**\newline%
% \newline%
% **Input:** words = \textbackslash{}{[} "cat ", "cats ", "catsdogcats ", "dog ", "dogcatsdog ", "hippopotamuses ", "rat ", "ratcatdogcat "\textbackslash{}{]}\newline%
% **Output:** \textbackslash{}{[} "catsdogcats ", "dogcatsdog ", "ratcatdogcat "\textbackslash{}{]}\newline%
% **Explanation:**  "catsdogcats " can be concatenated by  "cats ",  "dog " and  "cats "; \newline%
%  "dogcatsdog " can be concatenated by  "dog ",  "cats " and  "dog "; \newline%
%  "ratcatdogcat " can be concatenated by  "rat ",  "cat ",  "dog " and  "cat ".\newline%
% \newline%
% **Example 2:**\newline%
% \newline%
% **Input:** words = \textbackslash{}{[} "cat ", "dog ", "catdog "\textbackslash{}{]}\newline%
% **Output:** \textbackslash{}{[} "catdog "\textbackslash{}{]}\newline%
% \newline%
% **Constraints:**\newline%
% \newline%
% *   `1 <= words.length <= 104`\newline%
% *   `1 <= words{[}i{]}.length <= 30`\newline%
% *   `words{[}i{]}` consists of only lowercase English letters.\newline%
% *   All the strings of `words` are **unique**.\newline%
% *   `1 <= sum(words{[}i{]}.length) <= 105`'\newline%
% \}\}\newline%
% \newline%
% \{\{\newline%
%     "label": 1\newline%
% \}\}\newline%
% \newline%
Now, consider the question below and categorize it using the criteria above. Output your answer in a json format:\newline%
\{\{\newline%
    "title": "{\color{blue}Container With Most Water}",\newline%
    "question": "{\color{blue} You are given an integer array `height` of length `n`. There are `n` vertical lines drawn such that the two endpoints of the `ith` line are `(i, 0)` and `(i, height{[}i{]})`.\newline%
    \newline%
    Find two lines that together with the x{-}axis form a container, such that the container contains the most water.\newline%
    \newline%
    Return \_the maximum amount of water a container can store\_.\newline%
    % \newline%
    % **Notice** that you may not slant the container.\newline%
    % \newline%
    % **Example 1:**\newline%
    % \newline%
    % **Input:** height = \textbackslash{}{[}1,8,6,2,5,4,8,3,7\textbackslash{}{]}\newline%
    % **Output:** 49\newline%
    % **Explanation:** The above vertical lines are represented by array \textbackslash{}{[}1,8,6,2,5,4,8,3,7\textbackslash{}{]}. In this case, the max area of water (blue section) the container can contain is 49.\newline%
    % \newline%
    % **Example 2:**\newline%
    % \newline%
    % **Input:** height = \textbackslash{}{[}1,1\textbackslash{}{]}\newline%
    % **Output:** 1\newline%
    % \newline%
    % **Constraints:**\newline%
    % \newline%
    % *   `n == height.length`\newline%
    % *   `2 <= n <= 105`\newline%
    % *   `0 <= height{[}i{]} <= 104`}"\newline%
    \textit{rest of the question omitted...}"
}\\
\}\}\newline%
\\
{\color{violet}\{
    "label": 2
\}}
\\
       \bottomrule
    \end{tabular}
\end{table*}
